\subsection{The cover carrying the Gaussian norm tensor}

Let \(E:y^2=x^3-x\), let \(T_0=(0,0)\),
\(T_+=(1,0)\), \(T_-=(-1,0)\), and put
\[
                         F(Q)=\frac{x(Q)^2+1}{2y(Q)}.
\]
The notation \(\tau_T\) means translation by \(T\).

\begin{lemma}\label{lem:gaussian-norm-v4-cover}
The morphism \(F:E\to\mathbf P^1\) is a Galois cover
of degree four with deck group
\[
 \{1,\tau_{T_0},\tau_{T_+}[-1],\tau_{T_-}[-1]\}
                              \simeq(\mathbf Z/2)^2.
\]
Its branch values are \(1,-1,i,-i\).  It restricts to
an unramified cover
\[
 E\setminus E[4]\longrightarrow
       \mathbf P^1\setminus\{0,\infty,1,-1,i,-i\}.
\]
For every odd prime \(\ell\), the values \(F(Q)^4\),
with one \(Q\) in each order-four CM orbit on
\(E[\ell]\setminus\{0\}\), form a configuration of
\(M=(\ell^2-1)/4\) distinct points of
\(\mathbf P^1\setminus\{0,1,\infty\}\).
\end{lemma}
\begin{proof}
The addition formulas give
\[
 \tau_{T_0}(x,y)=(-1/x,y/x^2),\qquad
 \tau_{T_+}(x,y)=((x+1)/(x-1),-2y/(x-1)^2).
\]
Thus \(F\circ\tau_{T_0}=F\),
\(F\circ\tau_{T_+}=-F\), and \(F\circ[-1]=-F\).
The four displayed automorphisms preserve \(F\) and form
a four-group.  The poles of \(F\) are the four points
of \(E[2]\), each simple, so its degree is four.  This
proves the deck-group assertion.

The nontrivial translation has no fixed points.  The fixed
points of \(\tau_{T_+}[-1]\), respectively
\(\tau_{T_-}[-1]\), satisfy \(2Q=T_+\), respectively
\(2Q=T_-\).  Since \(F(Q)^2=x(2Q)\), their images are
respectively \(\{1,-1\}\) and \(\{i,-i\}\).
The fibers over zero and infinity are respectively
\([2]^{-1}(T_0)\) and \(E[2]\).  Together with the four
branch fibers, these are exactly \(E[4]\), proving the
unramified assertion.

The function \(F(Q)^4=x(2Q)^2\) is the order-four CM
quotient of \(2Q\).  Multiplication by two is a bijection
on \(E[\ell]\), commuting with the CM automorphism.
Thus distinct CM orbits give distinct values.  None of
these values is zero, one, or infinity, since the omitted
fibers consist of four-torsion points.
\end{proof}

This describes a second, explicit cover which carries the
product tensor.  Let
\(\mathcal C_M=\operatorname{Conf}_M
(\mathbf P^1\setminus\{0,1,\infty\})\) be the space of
unordered configurations of \(M\) distinct points.
The product of their affine coordinates is a regular unit
on \(\mathcal C_M\).  The equation
\begin{equation}\label{eq:gaussian-norm-configuration-cover}
                     c^4=\prod_{j=1}^{M}z_j
\end{equation}
therefore defines a finite etale cyclic cover of degree
four over \(\mathbf Q(i)\).  Its local monodromy is one
modulo four when one point travels around zero, minus one
around infinity, and zero around one or a collision
divisor.  These statements follow by the respective
orders of the product unit; in particular a braid
exchanging two points introduces no additional sign.

At the configuration in Lemma~\ref{lem:gaussian-norm-v4-cover},
the fourth-root fiber is exactly the norm tensor
\(C_\ell=\prod_Q F(Q)\).  The division-polynomial identity
identifies its fourth power as \(\chi_{-4}(\ell)/\ell\).
Consequently this is a specified finite cover and a
specified fiber in which to formulate the outstanding
full-GT comparison.

Preservation of the geometric monodromy class of
\eqref{eq:gaussian-norm-configuration-cover} does not yet
specify the action on its fiber at that configuration.
The missing assertion identifies the action transported
from the CM division configuration, including its
tangential normalization, with the Ihara Kummer action
on \(\ell^{-1/4}\).  Neither the four-group deck description
nor a normalizer statement for the full translation
group supplies that identification by itself.

One can also see the same missing descent in gamma
notation.  The adelic gamma function takes values in
Witt vectors over algebraic closures of the residue
fields, not necessarily in the finite residue field
containing a beta value; see
\cite[\S2.3, p.~101]{Furusho2003}.
An expression for a beta value as a square of such a
gamma value therefore does not prove that it is a square
in the rational norm-one torus.  The Frobenius action on
that square root is part of the remaining harmonic
residue condition.
